\documentclass[a4paper,11pt]{article}
\usepackage[utf8]{inputenc}
\usepackage[norsk]{babel}
\usepackage{amsmath}
\usepackage{graphicx}
\usepackage{geometry}
\usepackage{fancyhdr}
\usepackage{booktabs}
\usepackage{amssymb}

\geometry{margin=2.5cm}
\pagestyle{fancy}
\fancyhf{}
\lhead{Laboratorieoppgave: Reguleringsteknikk}
\rhead{Prosjekt: Den Digitale Tvillingen}
\cfoot{\thepage}

\begin{document}

\section*{Lab-oppgave: Systemidentifikasjon og SIMC-tuning}

\textbf{Gruppe:} \underline{\hspace{5cm}} \textbf{Dato:} \underline{\hspace{3cm}} \textbf{Anlegg nr:} \underline{\hspace{1cm}}

\subsection*{1. Identifikasjon av prosessen (Hvit og gul kurve)}
\textit{Kjør et sprang på tanken (Modus 0). OBS! Dere må ha is i magen og vente til den hvite kurven er helt flatlinjet før dere stopper målingen.}

\subsubsection*{Oppgaver MENS dere venter på sprangresponsen (10--15 min):}
\begin{enumerate}
    \item Tegn en grov skisse på frihånd av hvordan dere \textbf{tror} den hvite kurven vil se ut til slutt. Marker hvor dere antar at 63 \%-punktet vil lande på Y-aksen.
    \item Gå bort til det fysiske anlegget. Mål avstanden i centimeter fra pumpa til innløpet i tanken. Diskuter i gruppa hvordan denne rørlengden påvirker dødtiden ($L$) dere ser på skjermen.
    \item Slå opp SIMC-formelen for $K_p$. Hvis tidskonstanten ($T$) viser seg å bli veldig stor, vil regulatoren da bli aggressiv (høy $K_p$) eller forsiktig (lav $K_p$)? Begrunn svaret: \\
    \hrulefill
\end{enumerate}

\subsubsection*{Resultat av identifikasjonen:}
Når den gule kurven dekker den hvite perfekt, noter de endelige verdiene:
\begin{itemize}
    \item Forsterkning ($K$): \underline{\hspace{2cm}} 
    \item Tidskonstant ($T$): \underline{\hspace{2cm}} sekunder
    \item Dødtid ($L$): \underline{\hspace{2cm}} sekunder
\end{itemize}

\subsection*{2. Sertifisering for "Egendefinert Lambda"}
\textit{Før dere får lov til å bruke den manuelle inntastingsboksen i LabVIEW, må dere regne ut tålegrensen for anlegget deres manuelt og få godkjenning av lærer.}
\begin{enumerate}
    \item Regn ut målet for 5 \% oversving. Hvis setpunktet er 50 \%, hva er maksimalt tillatt nivå på tanken? \\
    \textbf{Maks nivå:} $50\,\% \times 1,05 = $ \underline{\hspace{3cm}}
    \item Bruk SIMC-formelen $K_p = \frac{1}{K} \cdot \frac{T}{\lambda + L}$. Hva skjer matematisk med forsterkningen ($K_p$) til regulatoren når dere tvinger $\lambda$ til å bli veldig liten? \\
    \hrulefill
\end{enumerate}
$\square$ \textbf{GODKJENT AV LÆRER (Signatur):} \underline{\hspace{4cm}} *(Åpner for manuell tasting)*

\newpage

\subsection*{3. Konkurranse: Finn den optimale Fartsfaktoren}
\textit{Bytt til regulatormodus (Modus 1). Målet er å oppnå høyest mulig fartsfart ($T/\lambda$) uten å passere grensen for 5 \% oversving på den blå kurven.}

\begin{table}[h]
\centering
\begin{tabular}{@{}lcccc@{}}
\toprule
\textbf{Strategi ($\lambda$)} & \textbf{Fartsfaktor ($T/\lambda$)} & \textbf{Integraltid ($T_i$)} & \textbf{Oversving (\%)} & \textbf{Stabil? (J/N)} \\ \midrule
$\lambda = T$                 & 1                                  &                              &                         &                        \\
$\lambda = T/2$               & 2                                  &                              &                         &                        \\
$\lambda = T/4$               & 4                                  &                              &                         &                        \\
$\lambda = T/6$               & 6                                  &                              &                         &                        \\
Egendefinert: \dots           &                                    &                              &                         &                        \\ \bottomrule
\end{tabular}
\end{table}

\subsection*{4. Analyse og diskusjon}
\begin{enumerate}
    \item Hva skjer med \textbf{Integraltiden ($T_i$)} når dere prøver å gjøre systemet raskere? \\
    \hrulefill
    \item Hvilken fartsfaktor ga den raskeste responsen uten å bryte 5 \%-regelen? \\
    \hrulefill
    \item Gruppe 2 fra forrige klasse opplevde at den blå kurven (virkeligheten) landet stabilt, men et par prosent \textbf{over} den rosa kurven (ønsket respons). Hvilken feil gjorde de i del 1 som førte til dette? \\
    \hrulefill
\end{enumerate}

\subsection*{5. Dokumentasjon}
\textit{Lagre PNG-bilde av deres beste forsøk. Sørg for at gruppenavn og tabellverdier er synlige i programmet før eksport.}

\end{document}
